\Needspace{18\baselineskip}
\section{Simultaneous substitution and shallow specimens}
\label{sec:simultaneous-specimens}

\Needspace{18\baselineskip}
\subsection{Simultaneous shallow substitution: carriers and incidence}
\label{sec:simultaneous}
The next milestone proves the multi-hole obstruction for a concrete
\emph{shallow} specimen representation. This needs a simultaneous filling
construction: the one-way side of a related pair may differ at different
holes, so one cannot simply assume that all left fillings, or all right
fillings, are one-way. The construction below works for arbitrary finite
typed graphs; only the later path arguments specialize to hypernets.

\begin{definition}[Simultaneous expansion]\label{term:simultaneous-expansion}
Let $C:I\to O$ be a finite typed graph and, for each edge $e$ of $C$, choose
a graph $H_e$ with exactly the source and target interface of $e$.
The expansion $C[H]$ replaces every edge at once. Its internal vertices and
edges have carriers
\[
 V_C+\sum_{e\in E_C}V_{H_e},\qquad \sum_{e\in E_C}E_{H_e}.
\]
An unchanged edge will be represented by its corolla. The vertex carrier
adds the same ordered external input and output constructors as before.
The host map $\iota_C$ retains each original vertex. For each $e$, the map
$\iota_e$ sends a filling internal vertex to its tagged fresh copy, its $i$th
input to $\iota_C(\operatorname{sourcePort}(e,i))$, and its $j$th output to
$\iota_C(\operatorname{targetPort}(e,j))$.
\end{definition}
\lean{Graph/Simultaneous}{12}{12}{Graph.Net.Expansion.Internal}
\lean{Graph/Simultaneous}{13}{13}{Graph.Net.Expansion.Vertex}
\lean{Graph/Simultaneous}{14}{14}{Graph.Net.Expansion.Edge}
\lean{Graph/Simultaneous}{16}{16}{Graph.Net.Expansion.host}
\lean{Graph/Simultaneous}{18}{21}{Graph.Net.Expansion.insert}

The host map is injective and its image is disjoint from every fresh internal
copy. Both facts follow by cases on the boundary/internal constructors and
on the disjoint-sum tag.
\lean{Graph/Simultaneous}{23}{24}{Graph.Net.Expansion.host_injective}
\lean{Graph/Simultaneous}{26}{27}{Graph.Net.Expansion.host_ne_internal}

\begin{lemma}[Uniqueness of a filling's owner]
If two non-output filling vertices have equal images, their owning host
edges are equal. The same holds for two non-input filling vertices.
Within one filling, the insertion map is injective on non-output vertices,
and separately on non-input vertices.
\end{lemma}
\begin{proof}
For the first assertion, two fresh internal images identify their owner tags.
An internal image cannot equal a boundary attachment in the host image.
The only remaining case compares two filling inputs. Their images are
source ports of their owning host edges; injectivity of the host map and
unique outgoing incidence identify those edges. For non-input vertices the
same argument uses output ports and unique incoming incidence.

Once the owner is fixed, two internal images identify their internal names;
two input images identify their input indices by distinctness of the host
source ports. The output argument uses target-port distinctness. No
injectivity is asserted between a filling input and output: a self-looping
host edge can identify them.
\end{proof}
\lean{Graph/Simultaneous}{29}{42}{Graph.Net.Expansion.source_owner}
\lean{Graph/Simultaneous}{44}{57}{Graph.Net.Expansion.target_owner}
\lean{Graph/Simultaneous}{59}{72}{Graph.Net.Expansion.insert_source_injective}
\lean{Graph/Simultaneous}{74}{87}{Graph.Net.Expansion.insert_target_injective}

\Needspace{18\baselineskip}
\subsection{Well-formedness of simultaneous expansion}
The source and target lists of a new edge $(e,x)$ are those of $x$ in $H_e$,
mapped by $\iota_e$. Each list remains duplicate-free: source vertices are
never outputs, target vertices are never inputs, and the respective
insertion injectivity lemmas therefore apply to every member of the list.
\lean{Graph/Simultaneous}{89}{89}{Graph.Net.Expansion.source}
\lean{Graph/Simultaneous}{90}{90}{Graph.Net.Expansion.target}
\lean{Graph/Simultaneous}{92}{94}{Graph.Net.Expansion.source_nodup}
\lean{Graph/Simultaneous}{96}{98}{Graph.Net.Expansion.target_nodup}

\begin{lemma}[Unique incidence in the expanded graph]
If two expanded edges share a source vertex, they are equal; the analogous
statement holds for a shared target vertex.
\end{lemma}
\begin{proof}
Pull the common incidence back along each mapped list. The owner lemma first
identifies the two original host edges. Insertion injectivity then identifies
the two filling vertices. Unique incidence within that one filling identifies
the two filling edges, hence the tagged expanded edges.
\end{proof}
\lean{Graph/Simultaneous}{100}{110}{Graph.Net.Expansion.source_unique}
\lean{Graph/Simultaneous}{112}{122}{Graph.Net.Expansion.target_unique}

\begin{lemma}[Boundary exclusion and coverage]
No expanded source is an output and no expanded target is an input. Every
vertex that is not an output has a source incidence, and every vertex that
is not an input has a target incidence.
\end{lemma}
\begin{proof}
For exclusion, inspect the local vertex witnessing a mapped incidence. A
fresh internal image has neither boundary status. A local input used as a
source maps to a host source port, which is not an output; a local output
cannot be a source. The target argument is dual.

For a retained host vertex requiring a source incidence, use host coverage
to find an original edge $e$ and its corresponding source-port index. Input
coverage in $H_e$ supplies an edge sourced at that filling input, whose image
is the original host vertex. Target coverage uses a filling output instead.
For a fresh internal vertex in a copy of $H_e$, apply coverage directly
inside $H_e$ and insert the resulting incidence. These cases exhaust the
expanded vertex datatype.
\end{proof}
\lean{Graph/Simultaneous}{124}{129}{Graph.Net.Expansion.source_not_output}
\lean{Graph/Simultaneous}{131}{136}{Graph.Net.Expansion.target_not_input}
\lean{Graph/Simultaneous}{138}{143}{Graph.Net.Expansion.source_host_covered}
\lean{Graph/Simultaneous}{145}{150}{Graph.Net.Expansion.target_host_covered}
\lean{Graph/Simultaneous}{152}{163}{Graph.Net.Expansion.source_coverage}
\lean{Graph/Simultaneous}{165}{176}{Graph.Net.Expansion.target_coverage}

Internal labels come from the indicated host or filling summand, and an edge
$(e,x)$ has the label of $x$ in $H_e$. The host map preserves vertex labels
by inspection. Each insertion map preserves them because its boundary cases
are exactly the host source-port and target-port typing equations.
\lean{Graph/Simultaneous}{178}{178}{Graph.Net.Expansion.internalLabel}
\lean{Graph/Simultaneous}{179}{179}{Graph.Net.Expansion.edgeLabel}
\lean{Graph/Simultaneous}{181}{182}{Graph.Net.Expansion.host_label}
\lean{Graph/Simultaneous}{184}{189}{Graph.Net.Expansion.insert_label}

Applying these label equations to each mapped incidence list reduces its
typing equation to the filling's original typing equation. Finite enumeration
is obtained by a dependent sum over the finite host-edge enumeration, and a
disjoint sum for the internal carriers. The constructor therefore proves all
fields of \texttt{Net}; it does not postulate the incidence invariants of the
simultaneous substitution.
\lean{Graph/Simultaneous}{191}{197}{Graph.Net.Expansion.source_typed}
\lean{Graph/Simultaneous}{199}{205}{Graph.Net.Expansion.target_typed}
\lean{Graph/Simultaneous}{211}{229}{Graph.Net.expand}

\Needspace{18\baselineskip}
\subsection{Retained corollas, identity expansion, and hole freedom}
\begin{definition}[Corolla copy]\label{term:corolla-copy}
A corolla copy of an original edge $e$ has one distinguished edge with the
same label as $e$, source list equal to the complete ordered input boundary,
and target list equal to the complete ordered output boundary. Every edge
of the copy equals the distinguished edge. The existing corolla constructor
supplies such a copy directly; its sole edge index is $0$.
\end{definition}
\lean{Graph/ExpansionPaths}{16}{21}{Linear.ExpansionPaths.CorollaCopy}
\lean{Graph/ExpansionPaths}{24}{30}{Linear.ExpansionPaths.corollaCopy}

Every source or target incidence of $e$ becomes the corresponding incidence
of its expanded corolla copy. Choose its port index in the host list; that
boundary vertex occurs in the corolla's full boundary list and maps to the
retained host vertex. More precisely, the copy's expanded incidence lists
are the original host incidence lists mapped by $\iota_C$. Mapping the full
port-index enumeration reconstructs the original ordered list, proving these
equalities. If the original edge label is globally focus-free, so is its
corolla copy: every edge of the copy is the distinguished one and has that
same label.
\lean{Graph/ExpansionPaths}{32}{37}{Linear.ExpansionPaths.copy_source}
\lean{Graph/ExpansionPaths}{39}{44}{Linear.ExpansionPaths.copy_target}
\lean{Graph/ExpansionFocus}{14}{19}{Linear.ExpansionPaths.copy_source_eq}
\lean{Graph/ExpansionFocus}{21}{26}{Linear.ExpansionPaths.copy_target_eq}
\lean{Graph/ExpansionFocus}{29}{33}{Linear.ExpansionPaths.copy_focusFree}

\begin{theorem}[Identity expansion]
Expanding every edge by its original corolla gives a graph isomorphic to $C$.
\end{theorem}
\begin{proof}
Forget the empty internal carrier of every corolla and the unique edge index
inside each corolla. These give explicit vertex and edge bijections; their
inverses retain a host vertex or attach edge index $0$. Internal and edge
labels are unchanged. For incidence, the exact copy-list equations reduce
the result to a host inclusion followed by its inverse projection, which
is the identity on every internal or boundary vertex.
\end{proof}
\lean{Graph/ExpansionUnits}{11}{19}{Linear.ExpansionPaths.identityVertices}
\lean{Graph/ExpansionUnits}{22}{31}{Linear.ExpansionPaths.identityEdges}
\lean{Graph/ExpansionUnits}{34}{65}{Linear.ExpansionPaths.identityIso}

\begin{lemma}[Hole-free expansion]
If every $H_e$ is globally hole-free, then $C[H]$ is globally hole-free.
\end{lemma}
\begin{proof}
A hole occurrence in the expansion lies in the label of some tagged edge
$(e,x)$, possibly within its nested box contents. That label is exactly the
label of $x$ in $H_e$. The same occurrence in that label would be a hole
occurrence in $H_e$, contradicting its global hole freedom.
\end{proof}
\lean{Graph/ExpansionUnits}{69}{74}{Linear.ExpansionPaths.holeFree}

\Needspace{18\baselineskip}
\subsection{Paired path reflection for many fillings}
A non-input vertex of $H_e$ cannot be inserted as a global input: target
coverage in $H_e$ supplies an incoming edge whose expanded copy still targets
its image. Global inputs have no incoming edge.
\lean{Graph/ExpansionPaths}{47}{50}{Linear.ExpansionPaths.insert_not_input}

\begin{theorem}[Local extraction under simultaneous filling]
If $\iota_e(v)$ is operation-reachable in $C[H]$, then $v$ is
operation-reachable in $H_e$, independently of all other fillings.
\end{theorem}
\begin{proof}
Induct on a reverse operation trace to the global input, generalizing $v$.
If $v$ is a filling input, take the empty local trace. Otherwise an empty
global trace is impossible by the previous lemma. In the nonempty case,
target coverage in $H_e$ provides an edge at $v$. Unique incoming incidence
in the expanded graph forces the first trace edge to be this inserted edge.
Pull its chosen source back into $H_e$, extract a local trace from the shorter
global trace by induction, and prepend that edge. This is the single-filling
extraction proof with the simultaneous incidence theorem providing uniqueness.
\end{proof}
\lean{Graph/ExpansionPaths}{54}{79}{Linear.ExpansionPaths.extract}

Consequently, if $H_e$ is one-way, no term-typed host target of $e$ can be
operation-reachable in $C[H]$. Write that target as a target port of $e$.
Its corresponding filling output has the same term type and the same image.
Extraction would make this output operation-reachable in $H_e$, contradicting
one-wayness.
\lean{Graph/ExpansionPaths}{83}{90}{Linear.ExpansionPaths.oneWay_target}

\begin{definition}[Paired output-closure for expansion]\label{term:paired-output-closure}
Two filling families $H,K$ satisfy the paired edge condition if, for every
original edge $e$, either both fillings are corolla copies of $e$, or $H_e$
is one-way, or $K_e$ is one-way. The disjunction is chosen separately for
each $e$; there need not be one side that is one-way at all replaced edges.
\end{definition}
\lean{Graph/ExpansionPaths}{97}{99}{Linear.ExpansionPaths.OutputClosedPair}

\begin{lemma}[Paired predecessor argument]
Suppose a term vertex $v$ of the host is the target of $e$ and its retained
copies are operation-reachable in both expansions. Under the paired edge
condition, $e$ is an operation, and every source of $e$ is a term vertex
whose copies are operation-reachable on both sides.
\end{lemma}
\begin{proof}
Either one-way branch contradicts reachability on that side by the previous
result. In the remaining branch both sides retain corolla copies. Their
expanded edges target the corresponding copies of $v$. The incoming-edge
theorem forces their label, hence the original label of $e$, to be an
operation. Every source of $e$ is therefore term-typed. Its mapped incidences
in the two copies and the predecessor reachability theorem give reachability
on each side.
\end{proof}
\lean{Graph/ExpansionPaths}{104}{122}{Linear.ExpansionPaths.predecessor}

\begin{theorem}[Paired host-path reflection]
If a host path ends at a term vertex reachable on both sides, all its edges
are operations. In addition, every source of every edge on the path is
term-typed and reachable on both sides.
\end{theorem}
\begin{proof}
Induct on the nonempty host path. For a singleton path use the predecessor
lemma. For a path $e::f::E$, apply the induction hypothesis to $f::E$. An
adjacency witness is both a target of $e$ and a source of $f$; the strengthened
induction conclusion supplies its type and reachability on both sides.
Apply the predecessor lemma to $e$, then combine its conclusion with that
for the tail. Repeated edges and any number of replaced holes are allowed.
\end{proof}
\lean{Graph/ExpansionPaths}{126}{150}{Linear.ExpansionPaths.path_operations}

\Needspace{18\baselineskip}
\subsection{Actual focussed expansions and their obstruction}
Retain the host focus by a corolla copy and require every other filling to
be globally focus-free. The expanded focus is the tagged distinguished edge
of that copy, with the same mode and the retained copies of the original
endpoints. The copy-list equalities give its singleton incidences. Host-map
injectivity preserves distinct endpoints, and the unchanged boundary
constructors preserve the prohibition on an exposed input-to-output focus.
For uniqueness, an edge from the focus copy equals its distinguished edge;
any edge from another filling is globally focus-free.
\lean{Graph/ExpansionFocus}{37}{55}{Linear.ExpansionPaths.focussed}

Initiality and finality are unchanged because both the mode and the input
status of the focus source are unchanged. Once global hole freedom of the
expanded graph is supplied, the construction packages the graph and its
proved focus as an actual state.
\lean{Graph/ExpansionFocus}{58}{62}{Linear.ExpansionPaths.focussed_boundary}
\lean{Graph/ExpansionFocus}{66}{71}{Linear.ExpansionPaths.state}

\begin{theorem}[Actual-state paired obstruction]
If the host has a path containing a non-operation edge and ending at its
focus source, and the two filling families satisfy the paired edge condition,
at least one of the two expanded states is not rooted.
\end{theorem}
\begin{proof}
If both states were rooted, Lemma~23 would make both actual focus sources
operation-reachable. They are the retained copies of the host focus source,
which has term type. Paired host-path reflection would therefore make every
edge of the given path an operation, contradicting its specified non-operation
edge. Classical case distinction turns this incompatibility into the stated
disjunction of non-rootedness.
\end{proof}
\lean{Graph/ExpansionFocus}{75}{93}{Linear.ExpansionPaths.pair_not_rooted}

\Needspace{18\baselineskip}
\subsection{Concrete pre-templates and shallow specimens}
\label{sec:shallow-specimens}
\begin{definition}[Concrete pre-template]\label{term:concrete-pre-template}
At a fixed depth bound, a concrete pre-template is an interface-indexed
binary relation on hypernets, with a proof that both members of each related
pair are globally focus-free. The relation and its domain condition are
defining data, not a law asserting operational behavior.
\end{definition}
\lean{Graph/ShallowSpecimens}{10}{12}{Linear.ConcretePreTemplate}

\begin{definition}[Concrete output-closure]\label{term:concrete-output-closure}
A concrete pre-template is output-closed when every related pair has a
one-way member on at least one side. This is the paper's output-closure
predicate, using the concrete one-way definition.
\end{definition}
\lean{Graph/ShallowSpecimens}{15}{16}{Linear.ConcretePreTemplate.OutputClosed}

For each original edge, the shallow specimen records one of two kinds of
data: retained corolla copies on both sides, or an actual hole label on
that edge together with a related pair of fillings. Thus arbitrary changes
to non-hole edges are not permitted by this specimen definition.
\lean{Graph/ShallowSpecimens}{20}{25}{Linear.SpecimenEdge}

\begin{definition}[Concrete shallow specimen]\label{term:shallow-specimen}
A shallow specimen consists of a globally unique-hole context of interface
$\star\to\epsilon$, a distinguished focus, and two simultaneous filling
families with the preceding edge-wise data. The erasure of the focussed
context must belong to the supplied allowed-context predicate. Both expanded
graphs must be globally hole-free. These fields express actual typing,
context, relation-membership, and closedness restrictions; they do not include
rootedness, path reflection, or a context-traversal law as assumptions.
\end{definition}
\lean{Graph/ShallowSpecimens}{30}{43}{Linear.ShallowSpecimen}

This is the shallow restriction of the paper's specimen construction. It
permits arbitrarily many outer holes and arbitrarily deep boxed contents in
the fillings. It does not replace holes inside the original host's boxes.
Compatibility with arbitrary addressed substitution and extension to deep
hole occurrences remain required before identifying it with the paper's full
specimen notion.

A focus label cannot equal a hole label, by disjointness of primitive labels
at either zero or successor depth. Therefore the focus edge must be in the
retained branch of the specimen data. This proves existence of its two
corolla copies, which are chosen to construct the focussed expansions.
\lean{Graph/ShallowSpecimens}{46}{48}{Linear.focusLabel_ne_hole}
\lean{Graph/ShallowSpecimens}{56}{61}{Linear.ShallowSpecimen.focus_copies}
\lean{Graph/ShallowSpecimens}{64}{64}{Linear.ShallowSpecimen.focusCopies}

Every other filling is globally focus-free. For a retained copy, use the
original focus's uniqueness and the copy-label theorem. For a filled hole,
use the pre-template's focus-free domain. These derived facts and the
specimen's proved graph closedness construct its two actual states.
\lean{Graph/ShallowSpecimens}{68}{73}{Linear.ShallowSpecimen.other_free}
\lean{Graph/ShallowSpecimens}{75}{76}{Linear.ShallowSpecimen.leftState}
\lean{Graph/ShallowSpecimens}{78}{79}{Linear.ShallowSpecimen.rightState}

\Needspace{18\baselineskip}
\subsection{Lemma 26(3) and the boundary bridge for shallow specimens}
\begin{theorem}[Output-closure supplies paired reflection]
An output-closed concrete pre-template supplies the paired edge condition
for every one of its shallow specimens.
\end{theorem}
\begin{proof}
At a retained edge, take the two recorded copies. At a filled hole, apply
output-closure to the recorded related pair to obtain a one-way side.
The side may change at the next hole.
\end{proof}
\lean{Graph/ShallowSpecimens}{83}{88}{Linear.ShallowSpecimen.outputClosedPair}

\begin{theorem}[Concrete shallow Lemma 26(3)]
For a shallow specimen of an output-closed pre-template, any context path
to its focus source containing a non-operation edge implies that at least
one of its two actual states is not rooted.
\end{theorem}
\begin{proof}
The preceding theorem supplies the paired edge condition. The focus-copy
and focus-freedom lemmas supply actual expanded states. Apply the actual-state
paired obstruction to those states and the given context path.
\end{proof}
\lean{Graph/ShallowSpecimens}{92}{100}{Linear.ShallowSpecimen.not_both_rooted}

Equivalently, when both specimen states are rooted, every edge of every
context path to the focus source is an operation. The formal proof applies
paired host-path reflection directly to the two operation-reachability
witnesses supplied by Lemma~23. This is the path consequence used in the
paper's Lemma~39(2), here proved for shallow specimens; it is not yet the full
entering/exiting classification or context-traversal theorem.
\lean{Graph/ShallowSpecimens}{104}{112}{Linear.ShallowSpecimen.rooted_context_path}

The proof uses the same incoming-incidence uniqueness and one-way extraction
as Proof~C.9. Rather than selecting a last hole and changing substitution
order, it reasons backwards simultaneously on both sides along the supplied
path. The paired predecessor lemma isolates exactly the contradiction that
the paper obtains at its last hole. This proof organization avoids assuming
an unproved interchange theorem.

\begin{theorem}[Concrete shallow Lemma 39(1)]
The two actual states of a shallow specimen agree on initiality and finality.
\end{theorem}
\begin{proof}
Each expanded focus has the same mode and source-input status as the original
context focus. Compose the left and right equivalences supplied by the
focussed-expansion boundary theorem. Output-closure and rootedness are not
needed for this boundary assertion.
\end{proof}
\lean{Graph/ShallowSpecimens}{116}{120}{Linear.ShallowSpecimen.boundary_agreement}

\Needspace{18\baselineskip}
\subsection{A constructed two-hole focussed specimen}
The earlier one-hole context has exactly one hole occurrence: inspection of
its labels rules out the focus and value edges, leaving its selected hole
edge. This proves global hole-label uniqueness and the absence of any
retained hole when that edge is removed.
\lean{Graph/ShallowSpecimenExamples}{10}{19}{Linear.Examples.oneWayContext_hole_only}
\lean{Graph/ShallowSpecimenExamples}{21}{23}{Linear.Examples.oneWayContext_unique}
\lean{Graph/ShallowSpecimenExamples}{25}{30}{Linear.Examples.oneWayContext_no_retained}

Replace that hole by the existing chain of two distinct holes. The exact
freshness theorem gives a globally unique-hole context, since the host
retains no hole and the inserted chain has distinct hole labels. Its focus
is retained by the previously checked shallow focus-free plugging operation.
\lean{Graph/ShallowSpecimenExamples}{33}{34}{Linear.Examples.twoHoleFocusedContext}
\lean{Graph/ShallowSpecimenExamples}{36}{38}{Linear.Examples.twoHoleFocusedContext_unique}
\lean{Graph/ShallowSpecimenExamples}{40}{41}{Linear.Examples.twoHoleFocus}

The simultaneous filling family replaces both new holes by variable corollas
and retains the focus and value by their original corollas. Define a
pre-template relating a graph to itself when it is focus-free and one-way.
This relation has the required focus-free domain, and is output-closed by its
one-way premise. At each hole the pair of variable corollas belongs to this
relation; at every other edge the corolla-copy fields hold by construction.
\lean{Graph/ShallowSpecimenExamples}{45}{50}{Linear.Examples.twoHoleFilling}
\lean{Graph/ShallowSpecimenExamples}{53}{55}{Linear.Examples.oneWayEquality}
\lean{Graph/ShallowSpecimenExamples}{57}{58}{Linear.Examples.oneWayEquality_closed}
\lean{Graph/ShallowSpecimenExamples}{60}{71}{Linear.Examples.twoHoleFilling_edges}

Every filling is globally hole-free. A variable corolla has no hole, while a
hole in a retained corolla would supply a retained hole occurrence in the
earlier one-hole context, already excluded. The expansion hole-freedom lemma
therefore proves closedness. These data assemble an actual two-hole specimen,
with the unrestricted allowed-context predicate for this example.
\lean{Graph/ShallowSpecimenExamples}{73}{82}{Linear.Examples.twoHoleFilling_closed}
\lean{Graph/ShallowSpecimenExamples}{85}{93}{Linear.Examples.twoHoleSpecimen}

The second hole by itself is a path to the focus source and is not an
operation. The concrete shallow Lemma~26(3) proves the required disjunction
of non-rootedness for the two constructed states, for every operation
classification. Their initial/final agreement follows from the concrete
specimen boundary theorem.
\lean{Graph/ShallowSpecimenExamples}{96}{105}{Linear.Examples.twoHoleSpecimen_not_rooted}
\lean{Graph/ShallowSpecimenExamples}{108}{111}{Linear.Examples.twoHoleSpecimen_boundary}

\Needspace{18\baselineskip}
\subsection{Alternating one-way choices: a necessary regression}
Use an operation signature with one unary operation, and form both its
operation corolla and a variable corolla. Relabel the existing two-hole chain
to this signature, leaving its hole interfaces and incidences unchanged.
On the left fill the first hole with a variable and the second with an
operation; on the right reverse these choices.
\lean{Graph/AlternatingFillings}{7}{7}{Linear.AlternatingFillings.ops}
\lean{Graph/AlternatingFillings}{8}{9}{Linear.AlternatingFillings.variableWire}
\lean{Graph/AlternatingFillings}{10}{11}{Linear.AlternatingFillings.operation}
\lean{Graph/AlternatingFillings}{14}{17}{Linear.AlternatingFillings.context}
\lean{Graph/AlternatingFillings}{19}{22}{Linear.AlternatingFillings.left}
\lean{Graph/AlternatingFillings}{24}{27}{Linear.AlternatingFillings.right}

The variable corolla is one-way: reachability of its output would force its
incoming variable edge to be an operation. The operation corolla is not
one-way, since its unique edge explicitly supplies an operation path between
its term input and output.
\lean{Graph/AlternatingFillings}{29}{34}{Linear.AlternatingFillings.variable_oneWay}
\lean{Graph/AlternatingFillings}{36}{38}{Linear.AlternatingFillings.operation_not_oneWay}

The paired edge condition holds: use the left one-way filling at the first
hole and the right one-way filling at the second. Neither entire family is
one-way edge by edge, as evaluation at the second left hole or the first
right hole yields the non-one-way operation corolla. Finally, the two-edge
host path cannot end at an output reachable on both sides. Otherwise paired
path reflection would make its hole labels operations, a contradiction.
This regression rules out an accidental strengthening of pointwise
output-closure to a single global one-way-side choice.
\lean{Graph/AlternatingFillings}{41}{51}{Linear.AlternatingFillings.output_closed}
\lean{Graph/AlternatingFillings}{54}{57}{Linear.AlternatingFillings.no_global_oneWay_side}
\lean{Graph/AlternatingFillings}{61}{72}{Linear.AlternatingFillings.not_both_reachable}
